\section{Introduction}

A player owes \$750 and holds \$168. Its first response proposes selling one Boardwalk house, mortgaging Park Place and \$335 of other deeds, and keeping Boardwalk and Kentucky unmortgaged. The recorded actions carry out that selection and leave \$28. Another player correctly values a much tighter estate at \$1,414, liquidates it to pay \$1,400, and retains the predicted \$14. These are different kinds of financial assessment. The first chooses resources under slack. The second values the full estate under a binding budget. Neither is adequately described by whether the final payment succeeds.

Economic agency also requires assessing other players' opportunities. A deed can complete the buyer's group while financing the seller's development. Refusing an offer can preserve a blocker while foregoing cash. Buying from a distressed rival can postpone an inheritance that would strengthen a third player. Monopoly provides observable versions of these relationships through fixed asset values, legal preconditions, persistent ownership, negotiated transfers, and uncertain future rents. Its financial arithmetic can be checked locally even when the strategic value of an exchange remains uncertain.

We ask what financial resources configured LLM agents represent in their first assessments, and how their applied actions and strategic communication relate to those assessments. The evidence consists of twenty complete MonopolyBench games involving four systems. We enumerate all explicit financing episodes and negotiation responses, then examine every first response in the sixteen episodes certified to require collateral unavailable in the initial mortgage menu. Contextual incidents connect economic statements to messages actually received, applied calls, and subsequent transactions. This is an observational study of an interacting roster, rather than a comparison of intrinsic model intelligence.

The clearest positive findings concern the contents of financial assessments. Agents explicitly describe sales enabling mortgages, produce ordered routes through those resources, calculate sufficient budgets, and sometimes select a subset of assets to preserve others. These observations coexist with precise distinctions between resource recognition, numerical correctness, and response consistency. Two failed maxima equal the same incomplete subtotal: cash, all building salvage, and initially available mortgages. Conversely, some successful responses contain numerical mistakes or unfinished provider fields. A third payable declaration contradicts its own stated rescue route.

The environment substantially assists execution. Reading its frozen source and the recorded portfolios establishes that continued valid cash raising suffices for every demonstrated-payable entry under specified conditions. No reference policy was executed for this paper. This structural result locates the contribution in the assessments and choices the agents express, rather than treating a long action sequence as proof of difficult planning. It also allows successful behavior to be described accurately without discounting the economic content of an explicit budget or a selective route.

The bargaining analysis adds the competitive setting in which these assessments matter. The complete response census reveals different ways of resisting an offer. Selected incidents show rival-development objections, correction of a false pitch, a public bargaining limit below private willingness, selective disclosure of finance, and a financial-denial appeal. A pooled completion contrast changes materially with proposer composition, so we retain it as an exploratory description. The paper's contribution is a source-linked account of financial assessment and economically situated conduct, with sufficient records to inspect both competence and its limits.

\section{Setting, records, and measurement}
\label{sec:setting}
\subsection{The observed economy and information}

Four players start with \$1,500 on the forty-space US board. They buy deeds, enter ascending auctions after declined purchases, trade cash and assets, and develop complete color groups. Bilateral negotiations allow counters and up to twenty exchanges per thread. Developed groups cannot be traded. Construction requires ownership, mortgage eligibility, even building, and bank stock. An explicit financing phase follows an obligation that cash cannot cover. It permits sales, mortgages, and voluntary declaration, refreshes the menu after each action, and automatically pays the current obligation when cash becomes sufficient. No trade or other player's move intervenes within that financing episode.

The mortgage-release rule matches the cited Hasbro rulebook: all buildings in the group must be removed before its deeds can be mortgaged~\cite{hasbroMonopolyRules}. Settlement differs. The frozen engine accepts voluntary declaration on a payable debt, transfers buildings intact to a player-creditor, charges no immediate interest on inherited mortgages, and does not immediately auction bank-returned assets. Its offered hotel downgrade requires four bank houses, whereas the cited rulebook also permits group hotel sales. These differences were not enumerated to agents. They affect the consequences of stopping and the interpretation of apparent rule knowledge. We distinguish a false financing budget from a strategic motive for surrender throughout.

The historical liquidation instruction says to choose one action, receive updated assets, and respond again if needed. It permits declaring when the player ``cannot or should not raise enough cash.'' The objective remains winning the whole game; it does not impose a separate task of paying every feasible debt. The standard prompt permits public deception and requests concise honest private notes. Accordingly, observed bluffing is conduct allowed by the task.

Each request supplies holdings, cash, mortgages, buildings, bank stock, player/model identities, twenty public timeline entries, and up to ten own notes. Negotiation requests include terms and thread history. Liquidation menus identify deeds by key without mortgage/salvage values or the release rule. Familiar names, earlier purchase information, carried estimates, and public actions remain available sources of knowledge. These are not memory-free tests of an unfamiliar transition system.

\subsection{Cohort and response conditions}

Five seeds each receive four cyclic seat rotations of one roster. Every system occupies each seat five times, while opponent order remains cyclic. Games last 75--386 turns, end without a turn cap, and contain 50,918 events, 6,989 resolved decisions, and 8,679 attempts. Shared seeds do not ensure identical later exposure because actions change random-draw consumption. Appendix~\ref{app:release} gives source commits, seeds, the game map, and replay provenance. All decision indices are zero-based.

\begin{table}[t]\centering
\begin{tabularx}{\linewidth}{@{}lYr@{}}\toprule
System & OpenRouter identifier & Decisions\\\midrule
Sol & \path{openai/gpt-6.1-sol} &1,961\\
Gemini & \path{google/gemini-4-argon} &1,669\\
Opus & \path{anthropic/claude-opus-5.5} &1,696\\
Grok & \path{x-ai/grok-4.7} &1,663\\\bottomrule
\end{tabularx}
\caption{Historical configurations. The client requests medium reasoning effort and does not set temperature, nucleus probability, or a token cap.}
\end{table}

Ordinary invalid outputs receive one corrective retry before deterministic fallback. A first refusal instead uses the remaining attempt for action-only recovery, which removes the private-note request. Opus records 1,604 first-attempt refusals, one additional later refusal, and no recovery refusals. All its first attempts in G5--G20 are refused. Grok uses required tool choice; the others use automatic choice. These unequal conditions define the evaluated configurations. They also affect opponents through offers, messages, and inherited state, so excluding one system's rows does not create an unaffected control.

We distinguish the applied tool call, its authored note, its public message, assistant content, and the provider reasoning field. The authored note is the tool argument historically named \texttt{private\_thought}; subsequent requests may carry it as memory. Provider reasoning is separately returned readable text, labeled \texttt{reasoning.summary} for Sol/Grok and \texttt{reasoning.text} for Gemini/Opus. It is not player memory or an opponent message, and may be incomplete or summarized. A valid received call is authoritative even if another field names a different action. The note and first action are jointly generated; the note is prospective relative to later collateral release, not an independently elicited forecast. JSON order provides no cognitive timing evidence.

\subsection{Populations and evidence criteria}

A financing episode begins at an explicit debt menu and ends in payment or elimination. Direct payments from sufficient cash are outside this population. Let $C$ be initial cash, $D$ debt, $S$ all nominal building salvage, $M_0$ initially offered mortgage principal, and $M_*$ all unmortgaged owned principal. Define
\[
 U_0=C+S+M_0,\qquad U_*=C+S+M_*.
\]
An upper bound below $D$ certifies insufficiency. Conditional on a demonstrated legal payment route, $U_0<D$ certifies that released mortgages are necessary. A sufficient nominal bound alone does not establish reachability. Feasibility labels combine observed payments, three retained historical payment witnesses, and one retained restricted-menu search. No engine execution, replay, policy run, search, or provider call was performed for this revision.

We code the first selected field as an explicit dependency, ordered route, inclusive budget, tentative resource, or omission. Selection means the authored note when present and otherwise the final attempt's provider field. Separate annotations describe numerical content and applied actions. These are model-assisted descriptive readings, without independent human validation or a claim about unexpressed cognition. A budget does not itself establish prerequisite knowledge; an unstated quantity is missing evidence, not an error. Both available channels for all sixteen cases are printed in Appendix~\ref{app:channels} and preserved with actions, memory, and provenance in the supplement.

An offer is a set of terms awaiting response. A counter changes its author. We classify group completion after the proposed transfer, including mortgaged deeds, and distinguish opening from later offers. Swaps transfer deeds both ways; cash bids request deeds for cash; cash asks offer deeds for cash. Counts refer to this corpus. Contrasts and incidents were selected retrospectively; resampling describes sensitivity to empirical weighting, not fresh replication or selection-adjusted inference.

\section{Financial assessment and financing choices}
\subsection{The complete financing census}

Of 188 financing episodes, 128 end in payment and sixty in elimination. Payments meet \$51,997 in obligations through 180 mortgage and 89 building-sale actions. Sixty-seven payments take one action; sixty-one take two to ten. No successful episode uses engine fallback. Three additional debts have retained legal payment witnesses despite declarations, giving 131 demonstrated-payable entries. Fifty-six eliminations exceed the generous bound. In the remaining G7 case, nominal resources are \$1,310 against \$750 owed, but the actual menu permits only \$440. Empty house stock blocks the hotel's downgrade and highest-level filtering hides the other pink houses. The retained search concerns the supplied menu, not every bundle accepted by the separate sale validator.

\begin{table}[t]\centering
\begin{tabularx}{\linewidth}{@{}Yrr@{}}\toprule
Resources required & Entries & Paid\\\midrule
Cash and initially offered mortgages &62&62\\
Sales needed; released mortgages unnecessary &53&52\\
Released mortgages necessary &16&14\\\midrule
Demonstrated payable &131&128\\\bottomrule
\end{tabularx}
\caption{Payable population. The middle row has 47 observed no-release routes, one retained sale-only witness, and five constructive source arguments. Its label does not follow from $U_0\geq D$ alone.}
\label{tab:partition}
\end{table}

The 62/53/16 partition in Table~\ref{tab:partition} separates resource requirements. For the middle 53, 47 actual payments use no mortgage absent initially, and the retained Grok witness uses only sales. Five further portfolios admit no-release routes by source reasoning: two contain houses only, two have stock for all hotel downgrades, and the last can downgrade one hotel and return its group's houses before the second downgrade. Appendix~\ref{app:proof} states those constructions. The sixteen necessary-collateral entries are 8.5\% of financing episodes and 12.2\% of payable entries. Nineteen actual payments use 33 newly released mortgages, including five where such use is not necessary.

\subsection{What continued financing guarantees}
\label{sec:continuation}

For all 131 demonstrated-payable entry states, uninterrupted valid, strictly cash-raising sales and mortgages suffice to pay the current obligation if the player does not voluntarily declare. This follows from source and entry holdings, rather than an executed baseline. Before payment, cash plus remaining salvage plus unmortgaged principal is invariant. Remaining building levels, treating a hotel as five, plus unmortgaged deeds strictly decrease. If assets cannot become inaccessible, exhaustion unpaid would leave cash equal to an initial value at least $D$, a contradiction.

All developed groups in these entries have complete ownership and legal evenness. In 108 hotel-free entries, a highest house remains sellable until the group clears. Twelve additional entries have sufficient initial mortgages, which remain available until used. Nine more have enough bank houses for every hotel downgrade even in the worst order. Two remaining G3 entries have one hotel per hotel-bearing group and recyclable stock. Their bank houses plus houses in hotel-free groups equal eight. Selling such a house preserves that sum; downgrading a group's sole hotel cannot decrease it. A hotel remains eligible when the bank has four houses; otherwise a hotel-free group supplies a sale. The argument therefore covers both remaining entries. Full premises are in Appendix~\ref{app:proof}.

The claim assumes valid positive-proceeds actions, uninterrupted execution, unchanged current debt, and no intervening actor. It excludes invalid-call loops, call/time limits, provider failures, and later obligations. It is neither a runtime guarantee for arbitrary agents nor a theorem about every nominally sufficient Monopoly state. Payment remains a recorded endpoint, but its occurrence here adds no evidence that a complete advance route was needed. Selective choices and expressed assessments can still contain information beyond that endpoint.

\subsection{Dependency statements, budgets, and selective routes}

The sixteen first fields divide into five explicit dependency statements, three ordered routes, five inclusive budgets, one tentative assessment, and two omissions. Nine successful cases use authored notes; five successful Opus cases use recovery provider fields; both failures use authored notes. The appendix gives row-level classifications and numerical annotations. The finer categories preserve what successful texts actually say without converting all resource inclusion into the same level of planning evidence.

G5 Gemini states that selling six dark-blue houses raises \$600, bringing \$74 cash to \$674, and releases \$175 and \$200 mortgages. The stated \$1,049 budget and \$149 post-payment remainder both match the records. Its provider field revises an earlier partial inventory within the same response. G18 d309 similarly describes selling eight pink houses to reach \$412, then mortgaging St. Charles for \$70 to pay \$450 and retain \$32. These statements identify the enabling relation before the mortgage menu exposes it.

G13 d195 supplies a complementary selective example. Gemini values \$445 of non-dark-blue mortgages, but chooses only \$335 of them together with the \$100 Boardwalk sale and \$175 Park Place mortgage. The planned balance is $168+100+335+175-750=28$. Its applied route mortgages Pacific, Electric, Indiana, and Park Place, while leaving Boardwalk and Kentucky unmortgaged as stated. With total nominal resources of \$1,088 against \$750 owed, this is a choice under slack. The record establishes selection and correspondence at the stated specificity, without identifying an optimum or a uniquely good asset-preservation policy.

Incident 1, G20 d327, instead illustrates exact aggregate valuation. Sol writes, ``Full liquidation should yield \$1,414 including cash and all mortgages, enough for \$1,400 rent by \$14.'' Two batched sale calls raise \$700, and eight mortgages raise \$665. With \$49 initial cash, payment leaves \$14. The final note names the last \$30 mortgage and the same remainder. Given these asset increments, payment requires the full-liquidation end state; action order is not unique. The ten calls document accounting and execution with feedback, not ten difficult route decisions. The creditor's same aggregate estimate contains offsetting component errors and is kept separate from the debtor's budget.

Success does not imply a fully accurate explanation. G2 recognizes the release but overstates yellow mortgages as \$430 rather than \$400. Its route still suffices. G20 d291 tentatively considers released pink mortgages while confusing the initial \$725 deficit with the post-financing shortfall. Cash \$275 plus \$450 salvage and \$230 initial mortgages gives \$955, leaving \$45, rather than the quoted \$725. It subsequently pays using released collateral. Other maximal budgets exceed what the actual route spends because auto-payment stops financing early. These are positive resource assessments with distinct numerical qualifications.

\subsection{Incomplete inventories and contrary channels}

Both focal failed budgets are correct additions over incomplete resources. In G12 d190, Sol gives $53+225+130=408$ as maximum finance against \$500, excluding \$340 released by clearing reds. Its note says no legal liquidation can save it. The provider field repeats the inventory and also considers that bankruptcy transfers property anyway. In G17 d234, Gemini gives \$500 salvage plus Baltic's \$30 mortgage as \$530 against \$550, excluding \$375 of dark-blue collateral. Its provider field ends by considering further options. The applied action declares in both cases; retained historical witnesses pay and leave \$38 and \$180 respectively.

Both failed totals equal $U_0$. Carried notes contain the same \$130 and \$30 mortgage figures, but those are also true current-menu values. Their recurrence does not identify memory as the source. Similar conditional inventories precede successful payments in G5, G13, G18 d313, and G20 d327. Menu emphasis, conditional-inventory reuse, rule application, output construction, and ordinary response variation remain possible explanations. The false maxima are established observations; the complete motives and internal causes of declaration are not.

The full population also contains six above-bound financing starts. Five are Grok calls whose other text names bankruptcy; one is Gemini mortgaging to impose an inherited-mortgage fee that the published rulebook requires but this engine does not charge. Table~\ref{tab:contrary} preserves each applied first action. G9's first call mortgages one deed; four mortgages occur across its episode. A separate payable Grok declaration, G12 d249, accompanies a note stating that bankruptcy would be a mistake and a sale route would pay. The structured declaration nevertheless executes.

\begin{table*}[t]\centering
\begin{tabularx}{\textwidth}{@{}llrrlY@{}}\toprule
Entry & System & Debt & $U_*$ & Applied first call & Other response evidence\\\midrule
G7 d286 & Grok &1,400&1,291&Mortgage Mediterranean&Content names declaration; provider field recognizes insufficiency.\\
G9 d194 & Grok &1,400&1,171&Mortgage Mediterranean&Note says declare now; content also names declaration.\\
G17 d396 & Grok &750&476&Sell light-blue houses&Content names declaration; note calls financing forced and insufficient.\\
G18 d344 & Gemini &400&382&Mortgage Virginia&Recognizes insufficiency and seeks to impose a 10\% inheritance fee absent from this engine.\\
G19 d224 & Grok &600&502&Mortgage Kentucky&Content names declaration; note incorrectly says reds must first be liquidated.\\
G20 d273 & Grok &450&417&Sell Illinois house&Note and content favor declaring immediately.\\\bottomrule
\end{tabularx}
\caption{All six financing starts above the generous bound. Continuing does not establish belief that payment is possible. None of these first calls is fallback.}
\label{tab:contrary}
\end{table*}

A narrow retained scan over 338 liquidation decisions finds nine conflicts between explicit action names in assistant content and applied calls. All are Grok, whose content names a parseable action in 27 of its 56 decisions. Other systems do not supply such named content here, so the scan cannot establish their absence of disagreement. Applied calls remain actions under the declared contract. We do not infer true intent from conflicting prose, attribute a resolver malfunction, or convert concordance with an asset bound into financial-decision accuracy.

\section{Bargaining over competitors' opportunities}
\subsection{A response census and its composition}

Across 526 threads, agents accept 113 offers, counter 394, and reject 413. Acceptance is legal in 918 of 920 opportunities; one counter and one rejection occur where acceptance is unavailable. Sol counters in 140/235 responses and Opus in 141/189; Gemini rejects in 180/242 and Grok in 149/254. Each accepts between 26 and thirty offers. Thus similar acceptance totals coexist with different ways of negotiating disagreement. Seven of Grok's 28 acceptances conflict with the action named in assistant content.

Opening swaps completing only the proposer receive 95/116 rejections, versus 37/72 completing neither participant, a 30.5-point difference. Excluding Grok-authored offers changes the counts to 90/108 versus 24/36 and the difference to 16.7 points. Grok authors half the neither-completing swaps and fifteen of their seventeen acceptances; thirteen of those fifteen originate in calls with conflicting content. Excluding conflicting openings instead gives 93/112 versus 26/43, a 22.6-point difference. These exclusions change the offer population and do not reconstruct corrected play.

Opening cash bids show 42/70 versus 67/103 rejections, a pooled difference of $-5.0$ points. Within-recipient differences range from $-0.2$ to $+8.5$ points. Later swaps show 8/39 versus 7/34, but are selected by earlier bargaining. Ownership completion also does not establish executable development: cash, mortgage status, stock, and opportunities to redeem matter. Compensation and opponent composition remain unmatched. Consequently these contrasts describe how offers were answered, rather than identifying an effect of completing a rival's group.

The appendix retains every form/completion cell and all five seed-block counts. Intact-block bootstrap enumeration gives a central no-Grok range of $[-3.8,31.9]$ points, while alternative game and two-stage reweightings also span zero. These are conditional sensitivity summaries over five observed blocks, not evidence of a null effect or independent confirmation. We report all raw small cells and avoid interpreting them as reliable system comparisons.

\subsection{What the bargaining language contributes}

A separate population comprises seventy-four legally acceptable recipient-completing rejections. Two earlier Astra readings agree that 61 mention a trade-linked consequence for the rival, including development or use of finance. Mere group naming is weaker evidence; a consequence statement need not be numerically correct. This inherited model-assisted count has no human validation, includes concentrated game/system exposure, and concerns offers disjoint from the proposer-only/neither comparison. It therefore cannot explain that quantitative contrast.

Incident 2, G15, shows a policy sustained across changing offers. At turn 63 Sol offers both greens for Ventnor and \$100, completing both groups. Grok refuses while contrasting Sol's roughly \$830 and financing capacity with its own \$233. That amount cannot fund a full \$600 green layer, although it can buy one house. Sol adds Kentucky and then \$250; Grok still refuses because it expects immediate yellow construction. Its turn-33 note already states the blocker policy, before a later public warning. The case documents persistent rival-focused resistance, rather than persuasion by that warning or proof that blocking improved winning chances.

Incident 3, G19, illustrates factual correction alongside acceptance. Grok pitches States as the third orange while offering Virginia and \$150. Both are pink. Sol replies, ``States is pink, not orange'' and accepts. The exchange completes no group. Grok's own note says to abort the mistaken proposal while its applied call sends it, so the sender's fields do not support one coherent deceptive plan. The recipient nevertheless distinguishes false justification from terms it is willing to accept.

Price histories add economic context. Seven seller/mortgage-status cells have higher median asks to completing buyers, but include completing-buyer samples as small as one. A different, closer comparison of six same-game/seller/deed histories with matched mortgage status has four higher, one lower, and one equal completing-buyer median. Fourteen earlier-resistance-to-settlement histories show falling holder cash in thirteen, but condition on eventual settlement and changing terms. These records support investigation of buyer identity and finance, without identifying a general monopoly premium or liquidity-caused concession.

\section{Strategic communication with financial stakes}

Incident 4, G7, connects an explicit bluff to a later, separate financing bargain. Grok publicly calls \$600 final for North Carolina while its private note says it would accept \$625. Sol settles at \$600. This is a documented difference between the public limit and stated private willingness under a prompt permitting deception. Later, Gemini observes Sol's \$84 cash after dark-blue construction and bids for New York. Sol asks \$700 to finance further construction while limiting the buyer's oranges; Gemini pays, mortgages five deeds, and builds four orange houses. Sol also adds dark-blue houses, then sells one to meet rent. The bargain finances competing investments. It is not the same transaction as the bluff, and changing terms prevent attributing settlement to cash alone.

Incident 5, G13, makes a rival's survival an object of negotiation. With \$35, Grok offers mortgaged Pennsylvania to Opus for \$200, arguing that otherwise Sol may inherit its yellows. Opus counters \$130. Grok's provider field explicitly plans to omit the fact that mortgaging Ventnor would make the lower payment sufficient. Opus pays \$200 plus \$16 transfer interest and considers both denying Sol the estate and potentially inheriting it itself. Its estate reasoning already appears after the opening appeal and before the omission-framed reply. Thus the records show selective presentation and uptake of estate considerations, without identifying the omission's effect on the extra \$70. Grok's incorrect 5/36 direct-roll probability is a separate factual error, not independently established deliberate probability deception.

Incident 6, G20, concerns financial denial. After Grok completes reds with little cash, Gemini publicly asks the others not to finance its houses by buying spare deeds. Grok's next request contains the appeal and its note recognizes it. Gemini later refuses Vermont for \$150, explicitly contrasting that house finance with a \$50 mortgage alternative. Sol had independently declined another purchase before the appeal. Opus later offers cash to fund red houses, although no sale settles. The evidence establishes an attempted coordination message and the speaker's follow-through; heterogeneous responses preclude a finding of collective compliance or a functioning cartel.

The full release preserves two further incidents in Appendix~\ref{app:incidents}: G2's buyer-specific yellow pricing and later reversal of ownership, including an incorrect development budget, and G18's public assurance with a simultaneous provider reservation followed by withdrawal. Together, the eight incidents cover competence, tactics, and inconsistencies. They are selected contextual examples, not prevalence estimates or evidence of tactics' causal effectiveness.

\section{Related work}

AgentMonopoly already studies LLM language, actions, and persona conditions in Monopoly~\cite{wu2026agentmonopoly}. Its reported negative-cash elimination differs from our explicit post-obligation financing phase. FinBoardBench studies several financial games with a common opponent pool and includes mortgages and building sales in Monopoly~\cite{hu2026finboardbench}. Multi-step financing is therefore not unique to this project. Whether its implementation permits declaration on payable debts is unverified. Our narrower object is the source-checkable connection between naturally encountered holdings, first financial assessments, voluntary terminal actions, and recorded execution.

Planning research already distinguishes applicable actions from reachable ones. ACPBench explicitly tests action reachability~\cite{kokel2024acp}; PlanBench and its o1 evaluation examine planning and false impossibility judgments~\cite{valmeekam2023planbench,valmeekam2024lrm}. AutoPlanBench compares plan generation with interactive action choice~\cite{stein2023autoplan}. ReAct studies interleaved action and feedback~\cite{yao2023react}. The Menu Is an Execution Prior includes a same-candidate ordering comparison~\cite{yan2026menu}. These precedents rule out treating resource release, false infeasibility, or generic menu sensitivity as discoveries here. Our familiar rule and refreshed legal menus warrant a financial-assessment interpretation, without a controlled claim about retrieval versus reasoning.

Statement/action studies examine declared-plan execution~\cite{oota2026declare}, while explanation-faithfulness work cautions against equating readable text with causal computation~\cite{turpin2023unfaithful,lanham2023faithfulness}. Memory research studies erroneous experience reuse~\cite{memory2026}; MemTrial examines financial experience dependence, with its PlantedMem construction using synthetic portfolio tilts rather than planting textual reserve figures~\cite{memtrial2026}. Our repeated reserve values do not establish a mechanism absent intervention. Saved-state probes and verifier assistance are established methods, exemplified by MafiaScope and LLM-Modulo~\cite{mafiascope2026,kambhampati2024modulo}; the present source proof is not a new agent architecture or executed policy comparison.

NegotiationArena studies controlled negotiation conditions~\cite{bianchi2024negotiationarena}; earlier dialogue agents separate strategic choices and language generation~\cite{lewis2017deal,he2018decoupling}. Identity-dependent externalities explain why a buyer's identity can matter to nonbuyers~\cite{jehiel1996externalities}, a useful interpretation of blocker and estate language without an equilibrium claim. Catan embargo simulations impose compliance and measure outcomes conditional on that restriction~\cite{guhe2014persuasion}; they do not establish autonomous persuasion. Vending-Bench and AutoVend use repeated operational trajectories and explicit references~\cite{backlund2025vendingbench,autovend2026}. The Arena communication study classifies a complete email population~\cite{li2026arena}. Our incidents instead preserve contextual acts and continuations, without claiming comparable population annotation or a causal conduct result.

\section{Discussion and conclusion}

The positive evidence is specific. Several agents represent resources absent from the initial menu, state correct monetary values, revise incomplete budgets, and choose selective financing. Bargaining records describe both participants' development budgets and the consequences of creditor succession. These observations would be lost in a payment fraction or an acceptance total. Conversely, the continuation proof, false maxima, and contrary channels prevent promoting expressed financial competence into a calibrated general planning score.

Persistent assets connect these local assessments to economic conduct. Trade-completed groups collect \$57,378 in the retained ownership histories; \$53,810 comes from groups built in the first management window. Thirty-eight of forty nonterminal eliminations follow developed-property rent. Estates transfer 227 deeds and complete 52 groups during play, plus 69 only at terminal elimination. These are historical stakes and exposure-dependent receipts, not causal investment returns. They explain why financing a rival or withholding a blocker can enter the agents' language without proving that the choices are advantageous.

Scope follows the design. Twenty games provide five seed blocks, one roster, no repeated seed/rotation cell, unequal response conditions, and retrospectively selected analyses. Model-assisted reading is disclosed once as such; multiple model reviews do not supply human validation. The records establish neither a causal memory or interface effect, optimal asset prices, durable learning, calibrated opponent models, nor a decomposition of winning into skill and luck. An external leaderboard would not by itself supply a matched capability axis. Appendix~\ref{app:questions} maps the original questions to supported observations.

The resulting account is an observational contribution about financial representations and strategic conduct under actual response conditions. Correct dependency statements, exact budgets, selective routes, false inventory maxima, and channel conflicts are separately inspectable. The distinction matters for evaluating interactive agents: successful payment can follow from continued valid action, while the financial content of the first response and the competitive consequences agents discuss require their own evidence.
